\chapter{局部 Langlands 参数的模空间}
\ysrc{3}
\begin{center}
{\small\itshape Jean-François Dat}
\end{center}
\medskip
\noindent\textbf{摘要.} 本文综述「局部 Langlands 参数的模空间」的动机、基础与主要结果，重点在于描述粗模空间，以及它们在一种「族中的」局部 Langlands 对应里预期扮演的角色。
\footnote{2020 年数学主题分类：主分类 11S37；次分类 11S20、14L30、22E50。}

\medskip
\noindent\textbf{目录.}
\begin{list}{}{\leftmargin=2.2em\labelwidth=2em\labelsep=0.4em\itemindent=-2.2em}
\item 动机：范畴化的 LLC 与「族中的」LLC
\item Langlands 参数。什么参数？
\item 如何处理连续性
\item 玩具模型：驯顺参数
\item 归约到驯顺情形与连通分支
\item 粗模空间及其应用
\item 参考文献
\end{list}

\section{动机：范畴化的 LLC 与「族中的」LLC}

局部 Langlands 纲领近来的若干引人注目的进展 [FS21] [Zhu25] [BZCHN20] [Hel23] 表明：用「伽罗瓦表示论」的数据来参数化 $p$-进群的不可约表示——这一最初由 Langlands 设想、随后经许多作者不断精确化的图景——只是一座冰山的一角，它应当在两侧关联远为一般的对象，而且是以更「范畴化」的方式关联。在参数一侧，这些更一般的对象就是局部 Langlands 参数的模空间上的（某种精细形式的）凝聚层。

于是这些模空间正在成为该领域的核心对象：局部 Langlands 对应预期的范畴化（CLLC）在参数一侧的舞台，正是在这里搭建的。不过，它们最早出现在 [Hel20b] 的 2016 年预印本版本中，那里只对 $\mathrm{GL}_n$ 的特殊情形作了引入。当时的动机是重新表述 [EH14] 中关于「族中的局部 Langlands 对应」（LLIF）之存在性的一个猜想。该猜想说的是：把 $\mathrm{GL}_n(F)$ 的表示的适当「族」指派给伽罗瓦群 $\Gamma_F = \mathrm{Gal}(\overline F/F)$ 的 $n$ 维表示，而系数取完备局部诺特环，其剩余域的特征 $\ell \neq p$。
\ysrc{4}该猜想的一个具体形态，是把伽罗瓦 $\mathbb F_\ell$-表示的通用形变环与整 Bernstein 中心在某些 $\mathbb F_\ell$-点处的完备化相比较。这样表述时，它基本上仍处在伽罗瓦表示形变理论这一成熟领域之内。然而 Helm 在 [Hel20b] 中的一个出色洞见，是把整个框架代数化。首先，他定义了 $\mathbb Z_\ell$ 上 $n$ 维 Weil 群表示的一个代数模空间，其各点的形式邻域恰好还原伽罗瓦表示论中通常的形变环。其次，他把 [EH14] 的猜想归约为：证明该模空间上的函数环与 $\mathrm{GL}_n(F)$ 上取值于 $\mathbb Z_\ell$ 的 Whittaker 函数空间的自同态环之间存在同构（在 $\mathrm{GL}_n$ 这一情形下，后者恰与整 Bernstein 中心同构）。这个同构最终在 [HM18] 中得到。

CLLC 与 LLIF 之间有一个先验的巨大差别：前者关键地用到精细模空间的叠性，而后者只通过函数环看到粗模空间。为了更精确而又非正式地一窥 CLLC 与 LLIF 的关系，我们先引入一些记号。

\begin{itemize}[leftmargin=1.6em,itemsep=0.2em]
\item $F$ 表示剩余特征为 $p$ 的局部非阿基米德域；
\item $\Gamma_F$ 表示「它的」伽罗瓦群，$W_F \subset \Gamma_F$ 表示其 Weil 群；
\item $G$ 表示 $F$ 上的约化群，$\hatG$ 表示其在 $\Spec \Z$ 上的对偶分裂约化群概形。由 $G$ 上的 $F$-有理结构，按 Langlands 的经典构造，我们得到 $W_F$ 在 $\hatG$ 上的一个保持固定数据（pinning）的作用；
\item $\Lambda$ 表示系数环。它主要是 $\Z_\ell$、$\Q_\ell$、$\mathbb F_\ell$ 或 $\Z[\frac1p]$ 之一，或是这些环的任一整扩张。无论如何，$p$ 在 $\Lambda$ 中总是可逆的；
\item $\Rep_\Lambda(G(F))$ 表示光滑 $\Lambda G(F)$-模的范畴。
\end{itemize}

回忆 Langlands 参数是 $W_F \to \hatG$ 的 1-上闭链，模去 $\hatG$-共轭（关于参数的更多内容见下文）。相应地，$G$ 的 Langlands 参数的模叠将具有形状 $Z^1/\hatG$，其中 $Z^1$ 是一个仿射 $\Lambda$-概形，参数化 1-上闭链，并带有 $\hatG$ 的自然作用。CLLC 的一项预测是：存在某个全忠实嵌入
\[
\Rep_\Lambda(G(F)) \hookrightarrow \Qcoh(Z^1/\hatG),
\]
把 $\Rep_\Lambda(G(F))$ 嵌入叠 $Z^1/\hatG$ 上拟凝聚层的（导出 $\infty$）范畴。此外，当 $G$ 拟分裂时，这样的嵌入应当只依赖于 Whittaker 数据 $(U,\psi)$ 的选取。更确切地说，联系于 $(U,\psi)$ 的嵌入应当把 Whittaker 表示 $W = W_{(U,\psi)} \in \Rep(G(F))$ 送到 $Z^1/\hatG$ 的结构层。作为全嵌入，这将特别蕴含同构
\[
\End(\mathcal O_{Z^1/\hatG}) = \Gamma(Z^1/\hatG, \mathcal O) = \Gamma(Z^1, \mathcal O)^{\hatG} \xrightarrow{\ \sim\ } \End(W).
\]
但这里的左端不过是粗模空间 $Z^1 \sslash \hatG$ 上的整体函数环，于是上述同构正是 LLIF 的主要预测——至少是 Helm 版本中的主要预测。
